\calculussheet{sheet:p08-inquisitive}{Natural deduction for \(\mathsf{InqB}\)}{Language: Definition~\ref{def:p08-inq-language}; \(\varphi,\psi,\chi\) arbitrary, \(\alpha\) syntactically classical. Negation and classical disjunction are defined. Inquisitive disjunction is \(\mathbin{\vee\!\vee}\). A displayed bracketed assumption is temporary; introduction/elimination may discharge any selected number of its occurrences, including zero. Other assumptions remain undischarged. No unrestricted substitution or unrestricted double-negation rule is added.}
\hbrule{rule:p08-inquisitive-and-i}{\(\land i\)}{\varphi\qquad\psi}{\varphi\land\psi}{\none}
\hbrule{rule:p08-inquisitive-and-e1}{\(\land e_1\)}{\varphi\land\psi}{\varphi}{\none}
\hbrule{rule:p08-inquisitive-and-e2}{\(\land e_2\)}{\varphi\land\psi}{\psi}{\none}
\hbrule{rule:p08-inquisitive-imp-i}{\(\to i\)}{[\varphi]^i\ \cdots\ \psi}{\varphi\to\psi}{Discharge the selected occurrences labelled \(i\) in the displayed subproof.}
\hbrule{rule:p08-inquisitive-imp-e}{\(\to e\)}{\varphi\qquad\varphi\to\psi}{\psi}{\none}
\hbrule{rule:p08-inquisitive-or-i1}{\(\mathbin{\vee\!\vee} i_1\)}{\varphi}{\varphi\mathbin{\vee\!\vee}\psi}{\none}
\hbrule{rule:p08-inquisitive-or-i2}{\(\mathbin{\vee\!\vee} i_2\)}{\psi}{\varphi\mathbin{\vee\!\vee}\psi}{\none}
\hbrule{rule:p08-inquisitive-or-e}{\(\mathbin{\vee\!\vee} e\)}{\varphi\mathbin{\vee\!\vee}\psi\qquad[\varphi]^i\ \cdots\ \chi\qquad[\psi]^j\ \cdots\ \chi}{\chi}{Discharge the selected \(i\)-occurrences in the first branch and \(j\)-occurrences in the second branch. The conclusion may be an arbitrary formula.}
\hbrule{rule:p08-inquisitive-bottom-e}{\(\bot e\)}{\bot}{\varphi}{\none}
\hbrule{rule:p08-inquisitive-split}{Split}{\alpha\to(\varphi\mathbin{\vee\!\vee}\psi)}{(\alpha\to\varphi)\mathbin{\vee\!\vee}(\alpha\to\psi)}{\(\alpha\) is classical; \(\varphi,\psi\) are arbitrary.}
\hbrule{rule:p08-inquisitive-dne}{\(\neg\neg e\)}{\neg\neg\alpha}{\alpha}{\(\alpha\) is classical.}
